\clearpage
\fancyhead[L]{\small TOBIAS MAYER. ASTRONOMISCHE BEOBACHTUNGEN.}
\hypertarget{native-649-lower}{}
\begin{center}\small Göttingische gelehrte Anzeigen. Stück 160. Seiten 1585..1588. 6. Oktober 1827.\\[3pt]\rule{.72\linewidth}{.3pt}\end{center}

\emph{Astronomical observations made at Göttingen from \textup{1756} to \textup{1761} by \textup{\textsc{Tobias Mayer}}. In two parts. Published by order of the commissioners of latitude.} London 1826. Die erste Abteilung umfasst 92 Seiten, die zweite 62 Seiten, beide im Folioformat.

Nach siebzig Jahren ist den Beobachtungen des großen Göttinger Astronomen \textsc{Tobias Mayer} noch die Auszeichnung zuteilgeworden, auf Kosten einer ausländischen Behörde gedruckt zu werden, die die großen ihr zur Verfügung stehenden Mittel gern zum Vorteil der Astronomie einsetzt. Die Handschrift der Beobachtungen aus den Jahren 1757..1761 war schon vor langer Zeit in den Besitz des Herrn von \textsc{Zach} gelangt, der sie der englischen Kommission für die Meereslänge mitgeteilt hatte. Die Handschrift der Beobachtungen des ersten Jahres befand sich noch im Archiv der hiesigen königl.\ Sozietät und wurde für den beabsichtigten Abdruck ausgeliehen.

Der bei Weitem größte Teil der Beobachtungen wurde am \textsc{Bird}'schen Mauerquadranten gemacht. \textsc{Mayer} hatte sein Hauptaugenmerk auf die Erstellung des Verzeichnisses der Zodiakalsterne gerichtet, und die Beobachtungen dieser Sterne bilden daher auch den Hauptinhalt des Tagebuchs. Sie beginnen am 8. Februar 1756. Zunächst finden sich bis zum 23. Mai bloß die Durchgangszeiten, und zwar bereits auf den mittleren Faden reduziert. Vom 23. Mai an erscheinen die Beobachtungen dagegen vollständig im Original. Die Beobachtungen des Jahres 1756 füllen die ganze erste Abteilung aus. Der Jahrgang 1757 nimmt in der zweiten Abteilung 35 Seiten ein, darauf folgen 6 Seiten bis in die Mitte des Jahres 1758. Hiermit enden \textsc{Mayers} Beobachtungen am Mauerquadranten eigentlich fast ganz. Am 5. April 1759 übergab er sein Verzeichnis der Zodiakalsterne der Königlichen Sozietät. Aus den ersten Tagen des Mai 1759 Beobachtungen des \textsc{Halley}'schen Kometen an einem parallaktisch aufgestellten Fernrohr. Aus dem Juni und Juli 1760 Beobachtungen am Mauerquadranten von \textsc{Mayers} würdigem Schüler \textsc{Niebuhr}. Schließlich vom 6. Juni 1761 die Beobachtung des Durchgangs der Venus durch die Sonne. Neben den Beobachtungen der Zodiakalsterne finden wir in den ersten Jahren noch sehr viele Beobachtungen der Sonne, mehrere der Venus und des Mondes und einige wenige des Mars, des Jupiters und des Saturns sowie verschiedene Sternbedeckungen.

Aus diesem Inhalt sieht man, dass eine von Grund auf neue Bearbeitung dieser Beobachtungen keine so umfassende Ausbeute liefern könnte, wie sie die gleichzeitigen Beobachtungen von \textsc{Bradley} durch \textsc{Bessels} vortreffliche Bearbeitung ergeben haben. \textsc{Mayer} hatte sich ein begrenztes \SourceLeaf{650}Hauptziel gesetzt, nämlich die Erstellung eines Verzeichnisses der Zodiakalsterne. Dieses hat er selbst auf seine Beobachtungen gegründet. Auch darf man nicht übersehen, dass die Dürftigkeit von \textsc{Mayer}'s Hilfsmitteln große Hindernisse dafür mit sich bringt, die Beobachtungen als absolute Beobachtungen zu reduzieren, und dass uns heute die Kenntnis einiger zusätzlicher Hilfsmittel fehlt, die er noch benutzte. \textsc{Mayer} hatte kein Mittagsfernrohr. Die Reduktion der am Mauerquadranten beobachteten Rektaszensionen setzt die Kenntnis der unregelmäßigen Abweichungen des Instruments vom Meridian voraus. Hierzu benutzte \textsc{Mayer} auch korrespondierende Sonnenhöhen an einem kleinen beweglichen Quadranten, doch sind uns solche Messungen nur von zwei Tagen erhalten. Ähnliche Schwierigkeiten gibt es bei den Zenitdistanzen. \textsc{Mayer} hatte keinen Zenitsektor. Eine zuverlässige Kenntnis des Kollimationsfehlers ließ sich daher nur durch Umhängen erlangen, und dieses Umhängen ist nur einmal geschehen. Alle vor dem Umhängen beobachteten Zenitdistanzen verlieren durch die von \textsc{Mayer} selbst gemachte Bemerkung, dass der Zentralzapfen, wie festgestellt wurde, locker saß, noch an Zuverlässigkeit. Es ließen sich noch mehrere andere Bemerkungen hinzufügen, wofür hier jedoch nicht der Ort ist.

Dagegen können noch oft einzelne Fälle eintreten, in denen es den Astronomen sehr willkommen sein wird, die Originalbeobachtungen nutzen zu können, wenn auch nur als Differenzialbeobachtungen. Bekanntlich hat sich die Handschrift schon einmal als sehr nützlich erwiesen, indem sie eine frühe Beobachtung des Uranus lieferte, den \textsc{Mayer} am 25. Sept. 1756 unwissentlich als Fixstern beobachtete.

In der Einleitung sind außerdem zwei Aufsätze aus den Opera inedita erneut abgedruckt, nämlich \textsc{Mayer}'s Vorlesung Observationes astronomicae quadrante murali habitae und \textsc{Lichtenberg}'s Bemerkungen über das von \textsc{Mayer} benutzte Thermometer. Am Schluss findet sich ein von Hrn.\ Prof.\ \textsc{Bessel} eingereichtes Druckfehlerverzeichnis zu \textsc{Bradley}'s Beobachtungen.
